\clearpage
\VerdierPhysicalPageAnchor{0202}%
\begingroup\small\setlength{\parskip}{0pt}\setlength{\emergencystretch}{1.5em}
\setlength{\abovedisplayskip}{4pt plus 2pt minus 2pt}%
\setlength{\belowdisplayskip}{4pt plus 2pt minus 2pt}%

\VerdierSemanticAnchor{ch3d-2-3-3-proof-q-r-step}
\[
\alpha\colon L^{-r+1}\longrightarrow Z^{-r+1},
\qquad\text{where }L^{-r+1}\text{ is an object of }M.
\]
Set $X_{q,r}=X_{q,r-1}(\alpha,-r+1)$ and
$v_{q,r}=p(\alpha,-r+1)$ (2.1.1). It is clear that $X_{q,r}$ and
$v_{q,r}$ have properties $I')_r$, $III')_r$, and
$IV')_r$. Moreover, by the definition of the complex $X_{q,r}$
(2.1.1), there is an exact sequence
\[
0\longrightarrow (X_{q,r})^{-r}\longrightarrow
L^{-r}\oplus L^{-r+1}\longrightarrow Z^{-r+1}\longrightarrow0.
\]
It then follows from property $(E_4)$ that $(X_{q,r})^{-r}$ belongs to
$M'$, and from (2.2.6) that
\[
M\text{-}\dim((X_{q,r})^{-r})
=M\text{-}\dim(Z^{-r+1})-1
=M\text{-}\dim(Y^q)-r-q,
\]
which is property $II')_r$.

\VerdierSemanticAnchor{ch3d-2-3-3-proof-q-completion}
We may therefore continue the construction. Set
$X_{q,-q+M\text{-}\dim(Y^q)}=X_q$, and let
$u_q\colon X_q\longrightarrow X_{q-1}$ be the morphism obtained by
composing the $v_{q,l}$. It is clear that $X_q$ and $u_q$ have properties
$1)_q$, $2)_q$, $3)_q$, $4)_q$, and $5)_q$. This completes the proof of
the first assertion of the proposition. To prove the second assertion, it
suffices to observe that, under the additional hypotheses, one may take
$X_{n_0}=Y$ in the preceding construction, for a suitable $n_0$.

\VerdierSemanticAnchor{ch3d-2-3-4}
\medskip\noindent\textbf{Remark 2.3.4.}\quad
In constructing the complex $X_{n_0}$ (cf.\ the proof of Proposition
(2.3.3)), we have in fact proved the following result. Let $M\subset M'$
be two sets of objects of $\mathcal A$ such that $0\in M$, the set $M'$
has properties $(E_1)$ and $(E_4)$, and every object of $M'$ is a quotient
of an object of $M$. Let
$Y\in\operatorname{Ob}\operatorname{comp}^{-}(\mathcal A)$ be a complex
all of whose components belong to $M'$. There is a complex
$X\in\operatorname{Ob}\operatorname{comp}^{-}(\mathcal A)$ all of whose
components belong to $M$, and a quasi-isomorphism
$u\colon X\longrightarrow Y$.

\VerdierSemanticAnchor{ch3d-2-3-5}
\medskip\noindent\textbf{2.3.5.}\quad
We leave it to the reader to state the results obtained by applying the
results of this section to the category opposite to $\mathcal A$. A set
$M$ of objects of $\mathcal A$ is said to have property $(E_i)^o$,
$1\leq i\leq6$, if the set $M$ of objects of $\mathcal A^o$ has property
$(E_i)$ of (2.2.1).

\VerdierSemanticAnchor{ch3d-2-4}
\setcounter{subsection}{3}%
\subsection{Applications.}

\VerdierSemanticAnchor{ch3d-2-4-1}
\noindent\textbf{Proposition 2.4.1.}\quad
\textit{a) Let $\mathcal A$ be an abelian category and $M$ a set of
objects of $\mathcal A$ having property $(E_1)$ (2.2.1). The category
$K^{-}(M)$ (1.1.4) is left localizing for
$\operatorname{Ac}(\mathcal A)$ (Chap. II, 2.3.5) if and only if the set
$M$ has properties $(E_2)$ and $(E_3)$.}

\noindent\textit{b) Let $M$ be a set of objects having properties
$(E_1)$, $(E_2)$, and $(E_3)$, and let $M^{\wedge}$ be the set associated
with $M$, defined in (2.2.1). The natural functor}
\[
K^{-}(M)/\operatorname{Ac}^{-}(M)
\longrightarrow
K^{-}(M^{\wedge})/\operatorname{Ac}^{-}(M^{\wedge})
\]
\endgroup
